% Qué datos se comparten con un asistente de IA: la señal no depende solo del color.
\documentclass[tikz,border=4pt]{standalone}
\usepackage{resonancia}
\begin{document}
\begin{tikzpicture}
  \foreach \i/\clase/\ej/\dec/\col in {0/Público/{Datos sintéticos del curso, normas, artículos abiertos}/{Sí}/rVerde, 1/Interno/{Borradores, apuntes, código en desarrollo}/{Con cautela}/rOcre, 2/Confidencial/{Registros de planta bajo acuerdo, resultados sin publicar}/{No, salvo permiso}/rNaranja, 3/{Datos personales}/{Listas de alumnado, notas, fotografías con personas}/{No, salvo base legal}/rNaranja, 4/Secretos/{Claves de API, contraseñas, tokens}/{Nunca}/rVioleta} {
    \pgfmathsetmacro{\y}{-\i*1.05}
    \draw[trazo, color=\col, fill=\col!8, rounded corners=3pt] (0,\y) rectangle (12.4,\y+0.9);
    \node[font=\sffamily\bfseries\small, anchor=west, color=rTinta] at (0.85,\y+0.45) {\clase};
    \node[etiqueta, anchor=west, color=rTinta, text width=5.0cm, font=\sffamily\footnotesize] at (3.9,\y+0.45) {\ej};
    \node[font=\sffamily\bfseries\small, anchor=east, color=\col!70!black] at (12.2,\y+0.45) {\dec};
  }
  % símbolos de cada fila (la señal no es solo el color): ✓ · ◐ · ✗ · ✗ · ✗✗
  \draw[rVerde, line width=2.2pt] (0.22,0.45) -- (0.38,0.25) -- (0.62,0.68);
  \draw[rOcre, line width=1.6pt] (0.42,-0.6) circle (0.2);  \fill[rOcre] (0.42,-0.6) ++(90:0.2) arc (90:270:0.2) -- cycle;
  \draw[rNaranja, line width=2.2pt] (0.22,-1.85) -- (0.62,-1.45) (0.22,-1.45) -- (0.62,-1.85);
  \draw[rNaranja, line width=2.2pt] (0.22,-2.9) -- (0.62,-2.5) (0.22,-2.5) -- (0.62,-2.9);
  \draw[rVioleta, line width=2pt] (0.12,-3.95) -- (0.38,-3.55) (0.12,-3.55) -- (0.38,-3.95) (0.45,-3.95) -- (0.71,-3.55) (0.45,-3.55) -- (0.71,-3.95);
\end{tikzpicture}
\end{document}
